\chapter{Introduction}
\label{chap:introduction}

\section{Research Context}
\label{sec:intro:context}

Detection systems built on machine learning are ordinary practice in security engineering.
Drebin classified Android applications from statically extracted feature sets in 2014
\cite{arp2014drebin}, MaMaDroid modelled sequences of abstracted API calls as Markov chains
soon afterwards \cite{onwuzurike2019mamadroid}, and a systematic review published in 2023
identified 132 studies applying deep learning to Android malware defence between 2014 and
2021 \cite{liu2023deeplearning}. The same shift took place in network defence, where anomaly
detection moved from statistical baselining \cite{chandola2009anomaly} towards deep
architectures trained directly on traffic \cite{pang2021deep}, and in program analysis, where
graph neural networks are now applied to code structure \cite{kipf2017semi,
velickovic2018graph, bilot2024survey}. Android alone reports more than three billion active
devices \cite{google2025android}, and the number of applications submitted for analysis each
year is far larger than any team of analysts can inspect by hand. That imbalance is what
makes automated classification attractive.

Industrial and Internet of Things (IoT) networks reached the same point by a different route. Their traffic is
comparatively regular, since the devices are few and their interactions are prescribed by the
process they control, which is what makes learned models of normal behaviour plausible there
at all. Testbeds built for the purpose supply the labelled attack data that operational plants
cannot \cite{mathur2016swat}. The difficulty is that a model of normal behaviour has to be
kept current as the deployment changes, and in a service-oriented architecture it changes
often.

The methods were developed for settings that differ from security in two ways that matter.
A classifier trained on one sample of data and deployed on another assumes that both come
from the same distribution. Malware violates the assumption by construction: families rise
and fall, developers repackage applications specifically to defeat the detectors of the
previous year, and the label an antivirus engine assigns to a file may change as engines are
updated. Jordaney et al.\ measured this drift directly and showed that a classifier's
confidence gives little warning of it \cite{jordaney2017transcend}. The second difference is
that the data are chosen by someone who benefits when classification fails.

How much reported performance survives a stricter protocol has been measured. Published F1
scores for Android malware classification reach 0.99, which leaves little room for
improvement and invites the question of what the number means. Pendlebury et al.\ traced the
figure to two forms of experimental bias: training and test data whose composition does not
match a real deployment, and time splits that let a classifier learn from samples that
appeared after the ones it is tested on \cite{pendlebury2019tesseract}. Arp et al.\ then
examined thirty papers from the four leading security conferences against a list of ten
methodological pitfalls. Every paper exhibited at least three of them \cite{arp2022dos}.
Artefacts in the datasets add to the effect. Drebin, the most widely used Android corpus,
carries duplicate samples produced by repackaging \cite{irolla2018duplication}, although the
systematic study of what those duplicates cost found the effect limited for supervised
classification and considerably larger for unsupervised family clustering
\cite{zhao2021impact}. None of these findings concerns a single weak paper. They concern the
standard protocol.

A second body of evidence concerns what happens once the adversary is treated as active.
Small, deliberately chosen perturbations flip the output of trained models
\cite{szegedy2014intriguing}, and the observation predates deep learning: Biggio et al.\
demonstrated gradient-based evasion against detectors at test time in 2013
\cite{biggio2013evasion}, with a line of work on adversarial classification running back
further still \cite{biggio2018wildpatterns}. The defensive response has not held up well.
Athalye et al.\ re-examined the nine non-certified defences accepted at ICLR 2018 and found
that seven owed their apparent robustness to obfuscated gradients, of which six were then
circumvented completely \cite{athalye2018obfuscated}. Tramèr et al.\ broke a further thirteen
defences published at leading machine-learning venues, all of which had already been
evaluated against adaptive attacks by their own authors \cite{tramer2020adaptive}. Carlini et
al.\ responded by writing down what an adequate robustness evaluation requires
\cite{carlini2019evaluating}. In the malware setting
the constraint is tighter than in image classification, because an adversarial example must
be a file that still installs and still runs \cite{pierazzi2020intriguing}.

Two questions follow from this, and they organise the present work. The first is a question
about representation: what should be observed about a program or a network, and how should
those observations be structured, so that a model is given the evidence that separates
malicious from benign behaviour? The second is a question about evaluation: how should a
defence be measured when the party generating the test data is trying to defeat it? The two
are usually pursued by different communities. This dissertation treats them together,
because a representation that is never tested against an adaptive opponent and a benchmark
that never varies the representation are weak in complementary ways.

\section{Problem Statement and Scope}
\label{sec:intro:problem}

\subsection*{The representation problem}

An analysis system sees only what its observation model admits. Static analysis reads the
application package and reasons about code that may never execute, which gives broad
coverage at the cost of accuracy against obfuscation and dynamic loading; the limits are
well understood and were formalised early \cite{moser2007limits}. Dynamic analysis observes
what a sample actually does, using instrumentation such as taint tracking
\cite{enck2014taintdroid}, which resists obfuscation but reaches only the paths that the
harness triggers \cite{onwuzurike2018family}. Neither observation model dominates the other,
so systems that combine them should be able to describe behaviour that neither describes
alone. What such a combination buys, measured on the same samples under one protocol, is the
first open question this dissertation addresses.

Structure raises a second problem. Program behaviour is commonly encoded as a graph, with
functions or basic blocks as nodes and calls or transfers as edges, and graph neural networks
then learn over that structure. An edge relates exactly two elements. Behaviour that is
security-relevant frequently involves more: a permission request, a sequence of API calls,
and a network write may be innocuous separately and malicious in combination. Encoding such a
group as a set of pairwise edges loses the fact that the members belong to one group.
Hypergraphs remove that restriction by allowing an edge to join any number of nodes
\cite{zhou2006learning}, and hypergraph neural networks provide the learning machinery
\cite{feng2019hypergraph, gao2023hgnn}. Whether the higher-order encoding
pays for itself on Android program structure, and whether weighting hyperedges by attention
adds anything beyond the encoding itself, is the second open question. The question is worth
asking sceptically, since attention mechanisms in graph learning have been shown to be less
expressive than their formulation suggests \cite{brody2022gatv2}.

\subsection*{The evaluation problem}

Robustness claims are only as good as the attacks used to test them, and the record shows
that attacks chosen by the defender's own authors are too weak
\cite{athalye2018obfuscated, carlini2019evaluating}. An evaluation that is credible has to
begin from an explicit statement of what the adversary knows and what the adversary may
change, and it has to be reproducible by a third party. Building such an evaluation as a
piece of infrastructure, rather than as a section of a paper, is the third open question
treated here.

\subsection*{Scope}

The empirical work concerns Android. Applications are analysed from the package as
distributed, using Dalvik bytecode and, where the design calls for it, traffic recorded while
the application runs. Native code, kernel-level instrumentation, and iOS lie outside the
scope. The industrial and IoT material concerns communication behaviour observed at the
network level and, in the SWaT study, sensor and actuator values of the physical process. It was
recorded in testbeds and project deployments, not at operator scale.

Three limits on the claims should be stated here rather than discovered later. First, the
work on communication graphs for IoT services was disseminated as a two-page contribution to
PAM 2018 on which the candidate is the fourth of four authors. The full paper that reports anomaly detection on
that system was written by two of his co-authors without him \cite{pahl2018all}, as was the
companion short paper on graph-based access control \cite{pahl2018graphbased}.
Chapter~\ref{chap:anomaly} identifies both and scopes the claim accordingly. Second, the
graph work on IoT services did not lead
to the Android work. The two are compared as representations in
Chapter~\ref{chap:discussion}, and no developmental sequence between them is claimed. Third,
the adversarial material is a contribution to evaluation methodology and to benchmark
construction. No new attack algorithm and no certified defence is proposed.

\section{Research Questions}
\label{sec:intro:rqs}

Four questions organise the dissertation. The first concerns the representation of behaviour
in networked and industrial systems and is divided into two subquestions. The second and third
concern the representation of Android program behaviour. The
fourth concerns evaluation. Each question takes up one of the four gaps stated in
Section~\ref{sec:background:summary}, and Chapter~\ref{chap:conclusion} answers each of them
in the same words in which they are asked here.

\researchquestion[1]{RQ1}{Behaviour in networked and industrial systems}{\rqOneText}

RQ1 takes up gap~\gapref{1}. Its two subquestions organise three empirical strands over service
relations, network records, and industrial process data. The strands are complementary rather
than cumulative, and no experiment compares all of them.

\researchquestion[1.1]{RQ1.1}{Communication graphs}{\rqOneOneText}

Chapter~\ref{chap:anomaly} answers RQ1.1 through a relation model for service-to-service
communication, a procedure for keeping the graph current as services change, and measurements
of the overhead this imposes. Detection accuracy against a labelled attack set is not
reported, and the chapter states why.

\researchquestion[1.2]{RQ1.2}{Modular industrial anomaly detection}{\rqOneTwoText}

Chapter~\ref{chap:anomaly} answers RQ1.2 through \nadics{}, a modular framework for network
and process features with interchangeable preprocessing and learning components. The chapter
reports its IUNO integration and a multivariate reconstruction detector on SWaT.
Here, reconstruction means that the model reproduces its input readings and uses the
discrepancies to help detect anomalies.

\researchquestion[2]{RQ2}{Multimodal Android malware analysis}{\rqTwoText}

RQ2 takes up gap~\gapref{2}. Chapter~\ref{chap:hybroid} answers it through \hybroid{}, which
extracts features from program structure and from recorded traffic and classifies their
combination. The chapter
reports binary detection and \categoryclassification{}, and it compares each modality against
the fusion under one protocol.

\researchquestion[3]{RQ3}{Higher-order Android malware analysis}{\rqThreeText}

RQ3 takes up gap~\gapref{3}. Chapter~\ref{chap:hgann} answers it through \hgann{}, which
represents an application as a hypergraph over its call structure and applies attention across
hyperedges. Two
classification tasks are distinguished throughout and are not interchangeable:
\familyclassification{} on Drebin, and \categoryclassification{} on CICMalDroid, in which
benign applications form one of the categories. Binary detection is a third, separate task.

\researchquestion[4]{RQ4}{Adversarial robustness evaluation}{\rqFourText}

RQ4 takes up gap~\gapref{4}. Chapter~\ref{chap:advml} answers it through a threat analysis of
machine-learning systems, a public contest run on a hidden test set, and a modular benchmark
in which attacks and defences are combined by configuration. It examines how task definitions,
scoring and attack validity affect the interpretation of the reported results, and identifies
the experimental records needed for reproduction.

\section{Contributions}
\label{sec:intro:contributions}

\paragraph{A communication-graph model for IoT service interaction.}
Service-to-service relations are represented as a graph that is built from observed traffic
and updated as the deployment changes, together with a measurement of the cost of maintaining
it. The contribution answers RQ1.1 and is reported in Chapter~\ref{chap:anomaly}. Its scope
is limited by the evidence available: the disseminated artefact runs to two pages, and the
chapter states what the candidate contributed to it.

\paragraph{A modular framework for anomaly detection in industrial networks.}
\nadics{} decomposes packet parsing, feature construction, and classification into
interchangeable components, so that learning methods can be exchanged without rebuilding the
pipeline, and it fed the visual security control room built by the \iuno{} project. The later
work adds a GAN-based detector for multivariate SWaT process data. The contribution answers
RQ1.2 and is reported in Chapter~\ref{chap:anomaly}.

\paragraph{Multimodal Android malware analysis.}
\hybroid{} combines a static representation of program structure with features derived from
observed network flows, and evaluates the combination against each modality alone on the same
samples. The contribution answers RQ2 and is reported in Chapter~\ref{chap:hybroid}.

\paragraph{Higher-order Android malware analysis.}
\hgann{} encodes groups of related program elements as hyperedges rather than as sets of
pairwise edges, and weights those hyperedges by attention. The contribution answers RQ3 and
is reported in Chapter~\ref{chap:hgann}, which compares the result against pairwise graph
models and against non-attentive hypergraph models. Because the encoding and the weighting
were changed together, the chapter does not attribute the gain to attention alone.
% T-012 answered by the author on 14 August 2026 and recorded as D-028: the ablation does not
% run. The isolation claim is withdrawn here and section 5.14 carries the confound as a
% declared limitation. Section 1.3 needs no change, because RQ3 compares systems rather than
% mechanisms.

\paragraph{Threat-guided evaluation of adversarial defences.}
An adversary model for machine-learning systems is translated into an evaluation protocol,
realised first as a contest with a hidden test set and then as a modular benchmark in which
attacks and defences are composed by configuration. The chapter derives the behaviour of
accuracy-loss scoring under constant predictions and distinguishes gradient masking from
prediction collapse. It analyses the response policy of binary adversarial detectors and
compares the published clean and attacked accuracies, with unresolved experimental mechanisms
identified as such.
The contribution answers RQ4 and is reported in Chapter~\ref{chap:advml}. Where the underlying
project deliverables were written collectively, the chapter states what the candidate
designed and implemented and what he synthesised from the work of partners.

\paragraph{A comparison across the contributions.}
Chapter~\ref{chap:discussion} relates the representations to one another along the axes of
observability, data-acquisition cost, and validity of evaluation. The comparison is
conceptual. Numerical results from the individual studies are not pooled, because their
datasets and validation protocols differ.

\section{Publications, Projects, and Individual Contributions}
\label{sec:intro:publications}

Parts of this dissertation have been disseminated previously, in accordance with the doctoral
regulations of the Technical University of Munich. Table~\ref{tab:intro:publications} lists
every such item, with the candidate's position in the author list and the chapter that draws
on it. Where a chapter reuses previously published material, it rewrites it rather than
reprinting it, and the corresponding results are restated from the original experimental
records.

\begin{table}[htbp]
  \caption{Previously disseminated parts of this dissertation.}
  \label{tab:intro:publications}
  \small
  \begin{tabularx}{\textwidth}{@{}Xlccc@{}}
    \toprule
    Item & Venue or type & Year & Position & Chapter \\
    \midrule
    Graph-based Anomaly Detection for IoT Microservices \cite{aubet2018graph}
      & PAM, two pages         & 2018 & 4 of 4      & \ref{chap:anomaly} \\
    \iuno{} work package 4, visual security control room \cite{iuno_ap4}
      & Project deliverable    & 2018 & contributor & \ref{chap:anomaly} \\
    \iuno{} work package 7, security tools \cite{iuno_ap7}
      & Project deliverable    & 2018 & contributor & \ref{chap:anomaly} \\
    \hybroid{} \cite{norouzian2021hybroid}
      & ISC, Springer LNCS     & 2021 & 1 of 4      & \ref{chap:hybroid} \\
    \hgann{} \cite{norouzian2025hgannmal}
      & ISC, Springer LNCS     & 2025 & 1 of 2      & \ref{chap:hgann} \\
    SPARTA D7.1, AI systems threat analysis \cite{sparta2021d71}
      & Public deliverable     & 2021 & contributor & \ref{chap:advml} \\
    SPARTA D7.5, AI systems security mechanisms \cite{sparta2021d75}
      & Public deliverable     & 2021 & contributor & \ref{chap:advml} \\
    SPARTA D7.6, Validation and evaluation report \cite{sparta2022d76}
      & Public deliverable     & 2022 & editor      & \ref{chap:advml} \\
    SAFAIR AI Contest development toolkit \cite{norouzian2021safairtoolkit}
      & Software repository    & 2021 & sole author & \ref{chap:advml} \\
    \bottomrule
  \end{tabularx}
\end{table}

% SETTLED 2026-08-14, T-007 and D-013. The two IUNO final deliverables are listed above and
% D4.2-1 is not. Reason: the text of AP4 and AP7 was published in the consortium results
% brochure, whereas D4.2-1 carries an internal copyright notice and no evidence of release.
% D4.2-1 is named in the individual-contribution prose below instead. Do not promote it into
% the table without evidence that it was disseminated, because this table is sworn to.

\subsection*{Individual contribution}

The two-page PAM 2018 contribution on graph-based anomaly detection for IoT microservices
\cite{aubet2018graph} arose from a collaboration in which the candidate is the last of
four authors. The peer-reviewed full paper on anomaly detection in that system was published
by two of those co-authors without him \cite{pahl2018all}, and so was the short paper on
graph-based access control in the same architecture \cite{pahl2018graphbased}.
Chapter~\ref{chap:anomaly} therefore presents the communication-graph material with its
authorship stated explicitly, and it limits the originality claim to the parts the candidate
designed and implemented.

Two final deliverables of the \iuno{} project \cite{iuno_ap4, iuno_ap7} report work the
candidate carried out at the chair. Both are consortium documents whose author lists sit at
the level of the document, so neither assigns any individual passage to one contributor.
Within the project the candidate designed and implemented \nadics{}, together with the
extractor that derives connection features from S7 and PROFINET traffic. He also produced the
labelled attack corpus the deliverables describe, recorded between Siemens S7-300 controllers
and their control server and labelled for binary detection as well as for five classes of
attack. An earlier internal deliverable of June 2017 \cite{iuno2017d42} documents the first
version of that tool and names him among its authors. That document was not published, which
is why it is absent from Table~\ref{tab:intro:publications}.
Chapter~\ref{chap:anomaly} states which components are his, and it records the passages in
which the deliverable text credits the implementation with functionality it did not yet have.

The candidate also designed, implemented, evaluated, and documented an unpublished extension
reported in Chapter~\ref{chap:anomaly}. It combines a multivariate GAN architecture with
reconstruction and critic scoring for SWaT process data. He carried out all of its experiments
and wrote the internal research report, which was not disseminated and is therefore absent from
Table~\ref{tab:intro:publications}.

% CONFIRM (T-007): the two sentences below describe individual contribution and will be
% sworn to in the Affidavit. Author position is verifiable from the papers; the division of
% design, implementation and experimentation is not.
% The \hybroid{} sentences are settled by D-025, answered by the author on 14 August 2026
% under Q1 and Q8 of T-032. Keep the partial qualification on the graph work: he made it
% himself and removing it would widen the claim. The \hgann{} sentence is settled by D-029,
% answered by the author on 14 August 2026 under Q11 of T-033.
For \hybroid{} \cite{norouzian2021hybroid} the candidate is the first of four authors, and
the paper records that the first two share first authorship. He designed the pipeline as a
whole and the network-traffic analysis within it, together with the framework through which
the dataset was analysed, from preprocessing through learning to the reported results. He
carried out the experiments and wrote the paper, and he contributed in part to the treatment
of the graph data. \hgann{} \cite{norouzian2025hgannmal} is the work of the candidate with
his supervisor. The candidate designed the method, implemented it, carried out the
experiments and wrote the paper.

Within the SPARTA project, the candidate contributed to the threat analysis of
machine-learning systems reported in \cite{sparta2021d71} and to the contest concept in
\cite{sparta2021d75}, and he is the editor of \cite{sparta2022d76}, in which the benchmark
and the contest results are reported. The first two deliverables were written by a
consortium and attribute contributors at the level of the document rather than the chapter.
Chapter~\ref{chap:advml} is therefore written afresh, cites each taxonomy it draws on in
place, and confines its contribution claim to the evaluation design and to the benchmark
implementation, which are the candidate's own.


\section{Dissertation Structure}
\label{sec:intro:structure}

Chapter~\ref{chap:background} sets out the foundations shared by the later chapters: Android
applications and their analysis, graph and hypergraph representations of programs, learning
over graphs, detection in networked and cyber-physical systems, adversarial machine learning,
and evaluation methodology. It closes with a statement of the gaps that the four research
questions address. The work against which each contribution is positioned is treated in the
chapter that reports it. Chapter~\ref{chap:anomaly} reports communication-graph, network-flow,
and industrial process-anomaly studies and answers RQ1.
Chapter~\ref{chap:hybroid} presents \hybroid{} and answers RQ2. Chapter~\ref{chap:hgann}
presents \hgann{} and answers RQ3. These two chapters carry the main Android empirical weight
of the dissertation. Chapter~\ref{chap:advml} presents the threat-guided evaluation work and answers
RQ4. Chapter~\ref{chap:discussion} sets the contributions against one another, and
Chapter~\ref{chap:conclusion} answers each research question in turn and states what remains
open.

Every contribution chapter carries its own account of the work specific to its method, its
own evaluation, and its own statement of limitations. There is no shared methodology chapter,
because the studies have no common experimental method to share.

Table~\ref{tab:intro:map} maps each research question to the chapter that answers it, to the
material that supports the answer, and to the candidate's role in producing that material.
The reader who wants to know what any particular claim rests on should start there.

\begin{table}[htbp]
  \caption{Research questions, chapters, and the material each answer rests on.}
  \label{tab:intro:map}
  \small
  \begin{tabularx}{\textwidth}{@{}llXl@{}}
    \toprule
    Question & Chapter & Supporting material & Role \\
    \midrule
    RQ1.1 & \ref{chap:anomaly}  & PAM 2018, 2~pp. \cite{aubet2018graph}, project artefacts & contributor \\
    RQ1.2 & \ref{chap:anomaly}  & \nadics{}, internal experimental records, \iuno{} deliverables \cite{iuno_ap4, iuno_ap7} & design, implementation, evaluation \\
    RQ2   & \ref{chap:hybroid}  & ISC 2021 \cite{norouzian2021hybroid}                            & first author \\
    RQ3   & \ref{chap:hgann}    & ISC 2025 \cite{norouzian2025hgannmal}                           & first author \\
    RQ4   & \ref{chap:advml}    & SPARTA D7.1, D7.5, D7.6 \cite{sparta2021d71, sparta2021d75, sparta2022d76} & contributor, editor \\
    \bottomrule
  \end{tabularx}
\end{table}
